\chapter{Conclusions}

\section{Research Outcomes}
Core contributions concentrated on probabilistic forecasting under extreme imbalance, asymmetric performance geometry, axiomatic contrasts among post-hoc explainers, and regulatory accountability lenses \cite{fraud_handbook,lime_original,shap_original,xai_survey,gdpr_text}. Empirical artefacts illustrate---without exhausting---those theoretical commitments.

Architecturally minimal instantiations reaffirm that explainability complements, never replaces, human adjudication whenever automated scores influence investigation queues \cite{fraud_handbook}. Methodological reproducibility dominates anecdotal toolchain preference.

\section{Answering the Research Questions}
\begin{description}
    \item[RQ1 (valuation under rarity)] Established metrics (precision--recall, threshold sweeps) supply principled substitutes for naive accuracy dominance; ROC summaries remain informative yet insufficient alone.
    \item[RQ2 (explainability semantics)] Cooperative-game attributions and local surrogate philosophies clarify admissible interpretive claims; hybrid masks $e_i(x)$ require explicit disclaimers distinguishing sensitivity from causal explanation.
    \item[RQ3 (probability artefacts)] Calibration and drift monitoring remain mandatory when adversaries exploit score surfaces; instantaneous probability outputs inherit dataset-specific ontology constraints.
\end{description}

Analyst-facing overlays that pair $\hat{y}$ with $e(x)$ improve transparency contingent on disciplined communication policies; versioning and reproducible artefacts operationalise reproducibility mandates \cite{xai_survey}.

\section{Practical Implications}
For financial institutions, the main practical implication is that explainable AI can be operationalized as a decision support capability rather than treated as a research-only add-on. A coherent implementation that joins prediction, explanation, and action logging can improve collaboration between data science teams, software engineers, and fraud operations analysts.

The work also highlights the importance of aligning model metrics with business constraints. In fraud detection, evaluation is not limited to aggregate accuracy; it must account for class imbalance, investigation capacity, and the cost of false negatives and false positives.

\section{Limitations}
Several limitations should be acknowledged. First, model behavior and explanation quality depend strongly on the representativeness and quality of the available transaction data. Second, local explanation methods provide useful evidence for individual predictions, but they do not establish causal relationships and may be sensitive to modeling assumptions \cite{shap_original}.

Third, operational performance claims must be interpreted in relation to the specific deployment setting, infrastructure, and analyst workflow in which the system is evaluated. Generalization to other institutions or fraud typologies requires additional validation.

\section{Future Work}
Future work can extend this dissertation along four directions. The first direction is broader evaluation across multiple datasets and temporal windows to assess robustness under distribution shifts. The second direction is integrating global explainability views and drift monitoring dashboards for model governance teams.

The third direction is introducing structured analyst feedback loops to support periodic retraining and policy-aware threshold adjustment. The fourth direction is comparative benchmarking of alternative explainability methods and model families under identical operational constraints, enabling evidence-based selection for production environments.

\section{Final Remarks}
Scientific novelty resides primarily in consolidating adversarial imbalance theory, disciplined interpretability taxonomy, and compliance-aware methodological design. Engineering supplements manifest these ideas but do not delimit scholarly boundary conditions. Subsequent scholarship should broaden distributional stress-tests, diversify explanation axioms under regulatory stress, and unify monitoring theory with behavioural studies of investigative trust.